\section{Experiments：Idea 生成与实验策略}
\label{sec:exp}

\webref{docs/4-实验/}
\resref{resources/04-experiment/}

\subsection{阶段规划}
粗略时间表为：想 idea 约 1--2 个月，做实验约 2--3 个月，写论文约 1 个月；实际往往更长。Idea-driven 方式跟进 SOTA 的 failure case，上手快但易撞车；Goal-driven 方式先设定 roadmap，更易形成体系但规划更难。

\subsection{从 Failure Case 到解法}
在更具挑战的数据设定上探索算法上限，将表面失败抽象为 technical challenge，再从可复用的 technique/insight「武器库」中组合解法。优秀工作往往不是简单模块堆叠，而是将要素巧妙结合。当领域出现新范式时，用其推进自身 roadmap 上的里程碑，常更易产生影响力。

\subsection{简单实验与记录}
简单实验应尽量只保留待验证的核心 challenge，排除无关干扰；验证有效后再迁移到真实数据并补齐对比与消融。实验记录需包含日期、步骤、想法、失败尝试、配置、结果与下一步。经典模型精读清单见 \path{resources/04-experiment/classic_models.md}，并配合李沐论文精读仓库推进。
